% =====================================================================
%  Self-Organising Warehouse Systems:
%  Graph-Based Coupling of Storage Location Assignment (SLAP)
%  and Multi-Agent Pickup and Delivery (MAPD) with RL Controllers
%
%  Notation follows the standard MAPF/MAPD convention G = (V, E)
%  (Ma et al., 2017; van Knippenberg et al., 2021; Pham & Bera, 2024)
% =====================================================================

% Suggested preamble packages:
% \usepackage{amsmath, amssymb, booktabs, longtable, array}

% ---------------------------------------------------------------------
\section{Introduction}
% ---------------------------------------------------------------------

Modern warehouses increasingly operate as tightly coupled sociotechnical
systems in which fleets of autonomous mobile robots and human co-workers
share the same physical infrastructure for order picking, replenishment and
internal transport. Two decisions dominate the efficiency of such systems.
The first is \emph{where items are placed}, studied as the Storage Location
Assignment Problem (SLAP) and, in its time-varying form, as the Dynamic
Storage Location Assignment Problem (DSLAP). The second is \emph{how agents
move and coordinate} in order to pick up and deliver items without
collisions, studied as Multi-Agent Path Finding (MAPF) and its lifelong,
online extension, Multi-Agent Pickup and Delivery (MAPD).

Both problems are naturally expressed on the same discrete structure. A
warehouse can be represented as a graph \( G = (V, E) \) whose vertices
\( V \) correspond to locations and whose edges \( E \) correspond to
connections between locations along which agents can move. In the MAPD
formulation, agents attend to a stream of delivery tasks, each characterised
by a pickup location and a delivery location in \( V \), and the quality of
a controller is measured by the average \emph{service time} needed to finish
a task after it entered the system. Storage assignment determines
\emph{which} vertices become pickup locations and how often, and therefore
directly shapes the spatial distribution of demand over \( V \) and \( E \).
Conversely, the routing policy determines which storage positions are
practically attractive, because a nominally close vertex is of little value
if the corridors leading to it are congested.

Despite this mutual dependence, the two problems are almost always studied
in isolation. Storage assignment is typically optimised under an assumed,
static routing cost such as a shortest-path or travel-time estimate, while
MAPF and MAPD are typically studied under a fixed layout and a fixed storage
policy. The resulting picture is incomplete: it cannot express the feedback
loop in which placement decisions change traffic patterns, and changed
traffic patterns in turn change which placements are good. Nor can it
express the failure mode that this feedback loop makes possible, namely that
individually rational decisions concentrate flow on a small subset of
aisles, storage vertices and paths, producing congestion and a commons-like
overuse of shared infrastructure.

This thesis studies exactly that feedback loop. The warehouse is modelled as
a single static graph \( G = (V, E) \). Storage assignment is modelled as a
time-varying map from items to storage vertices, updated from observed
transport and congestion outcomes. Routing is modelled as an online,
lifelong MAPD problem on the same graph, solved by controllers of three
different architectures: centralised, section-based, and decentralised
reinforcement-learning (RL) controllers that act on local projections of the
graph. The central question is whether closing the loop between storage
assignment and multi-agent routing improves collective behaviour, or whether
it destabilises it.

The choice of RL for the decentralised controller is deliberate rather than
incidental. Decentralised controllers must act under partial observability,
must react to congestion that they can only observe locally, and must do so
in real time as tasks arrive. RL provides a principled way to obtain such
reactive local policies without hand-designing a heuristic for every layout,
and recent work on graph-based MAPF shows that graph neural networks can
process arbitrary, non-lattice local neighbourhoods and support local
communication between nearby agents. This makes it possible to ask a
question that is not answerable with classical solvers alone: does giving
agents congestion-sensitive rewards and local communication change the
\emph{emergent} spatial organisation of the warehouse, and does it interact
constructively or destructively with a congestion-sensitive storage policy?

The contribution of the thesis is therefore not a new solver for either
subproblem. It is a graph-based formalisation of the coupled system, an
implementation of that system in simulation, and a controlled comparison
between decoupled and closed-loop configurations across controller
architectures, evaluated on travel distance per completed task, throughput,
congestion, flow concentration, and robustness under scaling in the number
of agents and the task arrival rate.

% ---------------------------------------------------------------------
\section{Notation}
% ---------------------------------------------------------------------

\begin{table}[htbp]
\centering
\caption{Notation used throughout the thesis. Graph, agent and task
notation follows the standard MAPF/MAPD convention.}
\label{tab:notation}
\small
\begin{tabular}{@{}ll@{}}
\toprule
Symbol & Meaning \\
\midrule
\multicolumn{2}{@{}l}{\textit{Warehouse graph (environment)}}\\
\( G = (V, E) \)            & Warehouse graph: vertices are locations, edges are connections \\
\( V \)                     & Finite set of vertices (locations) \\
\( E \subseteq V \times V \)& Set of edges; agents may move along an edge in one timestep \\
\( v, v' \in V \)           & Individual vertices \\
\( e = (v, v') \in E \)     & Individual edge \\
\( c(e) \geq 0 \)           & Traversal cost (time or distance) of edge \( e \) \\
\( \mathcal{N}(v) \)        & Neighbourhood of \( v \), i.e.\ \( \{ v' : (v,v') \in E \} \) \\
\( \mathcal{N}_d(v) \)      & Set of vertices within graph distance \( d \) of \( v \) \\
\( d_G(v, v') \)            & Shortest-path distance in \( G \) from \( v \) to \( v' \) \\
\( V_{\mathrm{str}} \subseteq V \)  & Storage vertices (shelves, bins, pallet positions) \\
\( V_{\mathrm{del}} \subseteq V \)  & Delivery vertices (pick stations, drop-off points) \\
\( V_{\mathrm{ep}} \subseteq V \)   & Endpoints: vertices at which agents may rest \\
\( V_{\mathrm{tsk}} \subseteq V_{\mathrm{ep}} \) & Task endpoints: possible pickup and delivery vertices \\
\( V_{\mathrm{ep}} \setminus V_{\mathrm{tsk}} \) & Non-task endpoints (parking / idle vertices) \\
\( \mathcal{Y} \)           & Partition of \( V \) into zones, \( V = \bigcup_{Y \in \mathcal{Y}} Y \) \\
\addlinespace
\multicolumn{2}{@{}l}{\textit{Agents}}\\
\( A = \{a_1, \dots, a_m\} \) & Set of agents; \( m = |A| \) \\
\( a_i \in A \)             & Agent \( i \) \\
\( t \in \{0, 1, \dots, T\} \) & Discrete timestep; \( T \) is the horizon \\
\( \ell_i(t) \in V \)       & Vertex occupied by agent \( a_i \) at timestep \( t \) \\
\( \ell_i(0) \)             & Initial vertex of agent \( a_i \) \\
\( \pi_i(t) \)              & Path of agent \( a_i \): mapping from timesteps to vertices \\
\( u_i(t) \)                & Action of agent \( a_i \) at timestep \( t \) (wait or move) \\
\( \mathcal{U}(v) \)        & Admissible actions at vertex \( v \) \\
\addlinespace
\multicolumn{2}{@{}l}{\textit{Items, orders and storage assignment (SLAP)}}\\
\( \mathcal{K} \)           & Set of item types (SKUs) \\
\( k \in \mathcal{K} \)     & An item type \\
\( \mathcal{Z} \)           & Order arrival distribution \\
\( z \in \mathcal{Z} \)     & An order, \( z = \{(k_1,q_1), \dots, (k_n,q_n)\} \) \\
\( q \in \mathbb{N} \)      & Requested quantity of an item type \\
\( \sigma_t : \mathcal{K} \to V_{\mathrm{str}} \) & Storage assignment at timestep \( t \) \\
\( \Sigma \)                & Class of admissible storage assignments \\
\( \rho_t(k) \)             & Estimated demand rate of item type \( k \) at timestep \( t \) \\
\( \mathrm{cap}(v) \)       & Storage capacity of vertex \( v \in V_{\mathrm{str}} \) \\
\( F \)                     & Storage update rule, \( \sigma_{t+1} = F(\sigma_t, \cdot) \) \\
\( \Delta \)                & Re-assignment interval (storage updated every \( \Delta \) timesteps) \\
\addlinespace
\multicolumn{2}{@{}l}{\textit{Tasks (MAPD)}}\\
\( \mathcal{T}(t) \)        & Task set: unexecuted tasks present at timestep \( t \) \\
\( \tau_j \in \mathcal{T}(t) \) & Task \( j \) \\
\( s_j \in V_{\mathrm{tsk}} \)  & Pickup vertex of task \( \tau_j \) \\
\( g_j \in V_{\mathrm{tsk}} \)  & Delivery vertex of task \( \tau_j \) \\
\( t^{\mathrm{arr}}_j \)    & Timestep at which task \( \tau_j \) enters \( \mathcal{T} \) \\
\( t^{\mathrm{fin}}_j \)    & Timestep at which task \( \tau_j \) is finished \\
\( \zeta_j \)               & Service time of task \( \tau_j \), \( t^{\mathrm{fin}}_j - t^{\mathrm{arr}}_j \) \\
\( f \)                     & Task arrival frequency (tasks added per timestep) \\
\addlinespace
\multicolumn{2}{@{}l}{\textit{Controllers and RL}}\\
\( \Pi \)                   & Class of routing controllers \\
\( \Pi_{\mathrm{cen}} \)    & Centralised controllers \\
\( \Pi_{\mathrm{sec}} \)    & Section-based (zone-decomposed) controllers \\
\( \Pi_{\mathrm{dec}} \)    & Decentralised RL controllers \\
\( \pi_\theta \)            & Parameterised policy with parameters \( \theta \) \\
\( o_i(t) \in \mathcal{O}_i \) & Local observation of agent \( a_i \) at timestep \( t \) \\
\( G_i(t) \)                & Local projection of \( G \) around \( \ell_i(t) \) \\
\( G^{A}(t) = (A, E^{A}(t)) \) & Agent interaction graph at timestep \( t \) \\
\( E^{A}(t) \)              & Interaction edges between agents within communication range \\
\( r \)                     & Interaction (communication) radius in \( G \) \\
\( x_i(t) \)                & Feature vector of agent \( a_i \) \\
\( h_i(t) \)                & GNN embedding of agent \( a_i \) after message passing \\
\( R_i(t) \)                & Reward of agent \( a_i \) at timestep \( t \) \\
\( \gamma \in (0,1) \)      & Discount factor \\
\( J(\pi_\theta) \)         & Expected discounted return of policy \( \pi_\theta \) \\
\addlinespace
\multicolumn{2}{@{}l}{\textit{Metrics}}\\
\( \bar{\zeta} \)           & Mean service time over all tasks \\
\( \Lambda \)               & Throughput: finished tasks per unit time \\
\( \bar{c} \)               & Mean travel cost per completed task \\
\( \mu(e) \)                & Traffic load on edge \( e \) (traversals per unit time) \\
\( \omega(e) \)             & Mean waiting time / queue length attributed to edge \( e \) \\
\( H \)                     & Flow-concentration index over \( E \) (e.g.\ Gini or entropy) \\
\( \mathcal{J} \)           & System-level cost functional of the coupled system \\
\bottomrule
\end{tabular}
\end{table}

% ---------------------------------------------------------------------
\section{Problem Statement}
% ---------------------------------------------------------------------

The thesis studies a coupled decision problem that is usually separated in
the literature. Storage location assignment is optimised under assumed
routing costs, while MAPF and MAPD are studied under a fixed layout and a
fixed storage policy, even though storage placement directly changes the
demand pattern over vertices and edges, and movement policies feed back into
which storage positions are practically attractive.

Informally, the problem is the following. A warehouse is given as a static
graph \( G = (V, E) \). Item types \( \mathcal{K} \) are stored at storage
vertices \( V_{\mathrm{str}} \) according to an assignment \( \sigma_t \).
Orders arrive from a distribution \( \mathcal{Z} \) and are decomposed into
delivery tasks, each with a pickup vertex determined by \( \sigma_t \) and a
delivery vertex in \( V_{\mathrm{del}} \). A fleet of \( m \) agents must
attend to this stream of tasks in an online, lifelong fashion while avoiding
collisions. Periodically, the storage assignment is revised using observed
transport and congestion outcomes, which changes the task distribution that
the routing controller faces. The question is under what conditions this
closed-loop coupling improves efficiency, congestion and robustness relative
to a decoupled baseline, and whether it induces commons-like concentration
of flow on a small subset of shared resources.

Formally, we require the following.

\paragraph{Given.}
\begin{itemize}
  \item a warehouse graph \( G = (V, E) \) with edge costs \( c : E \to \mathbb{R}_{\geq 0} \)
        and distinguished vertex subsets \( V_{\mathrm{str}} \), \( V_{\mathrm{del}} \),
        \( V_{\mathrm{ep}} \) and \( V_{\mathrm{tsk}} \);
  \item a set of item types \( \mathcal{K} \) with capacities \( \mathrm{cap}(v) \)
        for \( v \in V_{\mathrm{str}} \);
  \item an order arrival distribution \( \mathcal{Z} \) and task frequency \( f \);
  \item a set of agents \( A = \{a_1, \dots, a_m\} \) with initial vertices \( \ell_i(0) \);
  \item a class of admissible storage assignments \( \Sigma \) and an update
        interval \( \Delta \);
  \item controller classes \( \Pi_{\mathrm{cen}} \), \( \Pi_{\mathrm{sec}} \) and
        \( \Pi_{\mathrm{dec}} \).
\end{itemize}

\paragraph{Find.}
A pair \( (\sigma, \pi) \), consisting of a storage update rule
\( \sigma_{t+\Delta} = F(\sigma_t, \cdot) \) with \( \sigma_t \in \Sigma \) and a
routing controller \( \pi \in \Pi_{\mathrm{cen}} \cup \Pi_{\mathrm{sec}} \cup
\Pi_{\mathrm{dec}} \), such that the closed-loop system minimises the
system-level cost
\begin{equation}
  \mathcal{J}(\sigma, \pi)
  = \lambda_{\mathrm{trv}}\, \mathbb{E}\!\left[ \bar{c} \right]
  + \lambda_{\mathrm{svc}}\, \mathbb{E}\!\left[ \bar{\zeta} \right]
  + \lambda_{\mathrm{cng}}\, \mathbb{E}\!\left[ \textstyle\sum_{e \in E} \omega(e) \right]
  + \lambda_{\mathrm{cnc}}\, \mathbb{E}\!\left[ H \right],
  \label{eq:system-cost}
\end{equation}
where the expectation is taken over the order process \( \mathcal{Z} \) and any
randomness in the controller, and \( \lambda_{\bullet} \geq 0 \) are weights.

\paragraph{Compare.}
The closed-loop system is compared against the decoupled baseline in which
\( \sigma_t \equiv \sigma_0 \) for all \( t \), holding \( G \),
\( \mathcal{Z} \), \( f \), \( m \) and \( \pi \) fixed, and this comparison
is repeated across the three controller classes and across increasing
\( m \) and \( f \).

% ---------------------------------------------------------------------
\section{Formalisation}
% ---------------------------------------------------------------------

\subsection{Warehouse graph}

The warehouse is a graph
\begin{equation}
  G = (V, E),
\end{equation}
whose vertices \( V \) correspond to locations and whose edges
\( E \subseteq V \times V \) correspond to connections between locations
along which agents can move. Each edge \( e \in E \) carries a cost
\( c(e) \geq 0 \). The vertex set contains the distinguished subsets
\begin{equation}
  V_{\mathrm{str}} \subseteq V, \qquad
  V_{\mathrm{del}} \subseteq V, \qquad
  V_{\mathrm{tsk}} \subseteq V_{\mathrm{ep}} \subseteq V,
\end{equation}
where \( V_{\mathrm{str}} \) are storage vertices, \( V_{\mathrm{del}} \) are
delivery vertices, \( V_{\mathrm{ep}} \) are endpoints at which agents may
rest without blocking others, and \( V_{\mathrm{tsk}} = V_{\mathrm{str}} \cup
V_{\mathrm{del}} \) are the task endpoints. The graph is assumed connected,
and instances are assumed \emph{well-formed}: the number of non-task
endpoints is at least \( m \), and for any two endpoints there exists a path
between them traversing no other endpoint. Well-formedness is the standard
sufficient condition guaranteeing that lifelong instances remain solvable.

For hierarchical and section-based control we additionally fix a partition
\begin{equation}
  \mathcal{Y} = \{Y_1, \dots, Y_p\}, \qquad
  V = \bigcup_{n=1}^{p} Y_n, \qquad
  Y_n \cap Y_{n'} = \emptyset \ \ (n \neq n'),
\end{equation}
whose blocks correspond to aisles, sections or zones of the warehouse.

\subsection{Agents, actions and collisions}

Let \( A = \{a_1, \dots, a_m\} \). Agent \( a_i \) occupies vertex
\( \ell_i(t) \in V \) at timestep \( t \) and starts at \( \ell_i(0) \). At
each timestep an agent either waits or moves to an adjacent vertex,
\begin{equation}
  \ell_i(t+1) = \ell_i(t)
  \quad \text{or} \quad
  \big(\ell_i(t), \ell_i(t+1)\big) \in E ,
\end{equation}
so the admissible action set at vertex \( v \) is
\begin{equation}
  \mathcal{U}(v) = \{\textsf{wait}\} \cup \{ \textsf{move to } v' : v' \in \mathcal{N}(v) \}.
\end{equation}
Two collision constraints are imposed. Vertex collisions are forbidden,
\begin{equation}
  \ell_i(t) \neq \ell_{i'}(t)
  \qquad \forall\, a_i \neq a_{i'},\ \forall t,
  \label{eq:vertex-collision}
\end{equation}
and edge (swap) collisions are forbidden,
\begin{equation}
  \neg \big( \ell_i(t) = \ell_{i'}(t+1) \ \wedge\ \ell_{i'}(t) = \ell_i(t+1) \big)
  \qquad \forall\, a_i \neq a_{i'},\ \forall t.
  \label{eq:edge-collision}
\end{equation}
A path of agent \( a_i \) is a mapping from an interval of timesteps to
vertices; two paths collide iff the agents moving along them violate
\eqref{eq:vertex-collision} or \eqref{eq:edge-collision}.

\subsection{Storage assignment layer (SLAP)}

Let \( \mathcal{K} \) be the set of item types. A storage assignment at
timestep \( t \) is a map
\begin{equation}
  \sigma_t : \mathcal{K} \to V_{\mathrm{str}},
\end{equation}
subject to the capacity constraints
\begin{equation}
  \big| \{ k \in \mathcal{K} : \sigma_t(k) = v \} \big| \leq \mathrm{cap}(v)
  \qquad \forall v \in V_{\mathrm{str}} .
\end{equation}
The class of assignments satisfying these constraints is denoted \( \Sigma \).

Orders arrive from \( \mathcal{Z} \). An order is a finite multiset
\begin{equation}
  z = \{ (k_1, q_1), \dots, (k_n, q_n) \}, \qquad k_j \in \mathcal{K},\ q_j \in \mathbb{N},
\end{equation}
and each order line \( (k_j, q_j) \) induces one or more delivery tasks whose
pickup vertex is \( \sigma_t(k_j) \) and whose delivery vertex lies in
\( V_{\mathrm{del}} \). The storage assignment therefore determines the
spatial distribution of pickup vertices, and hence the demand induced on
\( V \) and \( E \).

Let \( \mu_t(e) \) denote the observed traffic load on edge \( e \) and
\( \omega_t(e) \) the observed waiting time attributed to \( e \), both
estimated from agent trajectories over the window \( [t-\Delta, t] \), and let
\( \rho_t(k) \) denote the estimated demand rate of item type \( k \). The
closed-loop storage update is
\begin{equation}
  \sigma_{t+\Delta}
  = F\big( \sigma_t,\ \rho_t,\ \mu_t,\ \omega_t \big),
  \qquad \sigma_{t+\Delta} \in \Sigma ,
  \label{eq:storage-update}
\end{equation}
where \( F \) trades off proximity against congestion. A concrete
instantiation assigns to each candidate storage vertex the score
\begin{equation}
  \phi_t(k, v)
  = \rho_t(k) \cdot \Big[
      \alpha \, \bar{d}_G\big(v, V_{\mathrm{del}}\big)
      + \beta \, \psi_t(v)
    \Big],
  \qquad
  \sigma_{t+\Delta}(k) \in \arg\min_{v} \phi_t(k, v),
  \label{eq:storage-score}
\end{equation}
where \( \bar{d}_G(v, V_{\mathrm{del}}) \) is the mean shortest-path distance
from \( v \) to the delivery vertices, \( \psi_t(v) \) aggregates congestion
on the edges leading to \( v \), and \( \alpha, \beta \geq 0 \) are weights.
Setting \( \beta = 0 \) recovers a purely distance-driven (congestion-blind)
storage heuristic; setting \( \Delta = \infty \) recovers the decoupled
baseline.

\subsection{Routing layer (MAPD)}

At each timestep the system adds all newly arrived tasks to the task set
\( \mathcal{T}(t) \). Each task
\begin{equation}
  \tau_j = (s_j, g_j), \qquad s_j, g_j \in V_{\mathrm{tsk}} ,
\end{equation}
is characterised by a pickup vertex \( s_j \) and a delivery vertex
\( g_j \). An agent is \emph{free} if it is not currently executing a task
and \emph{occupied} otherwise. A free agent may be assigned any task
\( \tau_j \in \mathcal{T}(t) \); it must then move from its current vertex via
\( s_j \) to \( g_j \). Upon reaching \( s_j \) it starts executing the task,
which is removed from \( \mathcal{T} \); upon reaching \( g_j \) it finishes
the task and becomes free again.

The service time of task \( \tau_j \) is
\begin{equation}
  \zeta_j = t^{\mathrm{fin}}_j - t^{\mathrm{arr}}_j ,
\end{equation}
and the primary routing objective is the mean service time
\begin{equation}
  \bar{\zeta} = \frac{1}{|\mathcal{T}_{\mathrm{fin}}|}
    \sum_{\tau_j \in \mathcal{T}_{\mathrm{fin}}} \zeta_j ,
\end{equation}
with throughput \( \Lambda = |\mathcal{T}_{\mathrm{fin}}| / T \) and mean travel
cost per completed task
\begin{equation}
  \bar{c} = \frac{1}{|\mathcal{T}_{\mathrm{fin}}|}
    \sum_{i=1}^{m} \sum_{t=0}^{T-1}
      c\big(\ell_i(t), \ell_i(t+1)\big).
\end{equation}

\subsection{Controller classes}

Three controller classes are considered on the same graph \( G \).

\paragraph{Centralised (\( \Pi_{\mathrm{cen}} \)).}
A single planner assigns endpoints to all agents and plans collision-free
paths for all agents jointly, re-planning at every timestep. This yields
strong global coordination at a computational cost that grows steeply with
\( m \).

\paragraph{Section-based (\( \Pi_{\mathrm{sec}} \)).}
Assignment and planning are decomposed over the zones \( \mathcal{Y} \).
Within a zone \( Y_n \) a local planner operates as in the centralised case;
coordination between zones is restricted to boundary vertices
\( \partial Y_n = \{ v \in Y_n : \mathcal{N}(v) \setminus Y_n \neq \emptyset \} \).

\paragraph{Decentralised RL (\( \Pi_{\mathrm{dec}} \)).}
Each agent selects its own action from a local observation using a shared
parameterised policy \( \pi_\theta \).

\subsection{Local projections and the agent interaction graph}

For the decentralised controller, agent \( a_i \) does not observe \( G \)
globally. Define the depth-\( d \) local projection of \( G \) at vertex
\( v \) as the induced subgraph
\begin{equation}
  G_i(t) = \mathrm{LocProj}\big(G, \ell_i(t), d\big)
         = G\big[\, \mathcal{N}_d(\ell_i(t)) \,\big],
  \qquad
  \mathcal{N}_d(v) = \{ v' \in V : d_G(v, v') \leq d \}.
  \label{eq:locproj}
\end{equation}
Because a goal vertex may lie outside \( G_i(t) \), each vertex
\( v' \in \mathcal{N}_d(\ell_i(t)) \) is annotated with the scalar potential
\begin{equation}
  \eta_i(v') = d_G\big(v', g_{j(i)}\big),
  \label{eq:potential}
\end{equation}
where \( g_{j(i)} \) is the current target of agent \( a_i \). These
potentials replace the Euclidean direction vector available on a grid:
progress towards the goal corresponds to selecting a neighbour with lower
\( \eta_i \), and \eqref{eq:potential} is precomputed once per instance.

Local interaction between agents is captured by the time-varying agent
interaction graph
\begin{equation}
  G^{A}(t) = \big(A, E^{A}(t)\big),
  \qquad
  (a_i, a_{i'}) \in E^{A}(t) \iff d_G\big(\ell_i(t), \ell_{i'}(t)\big) \leq r ,
  \label{eq:agent-graph}
\end{equation}
with communication radius \( r \). Note the deliberate distinction between
the static environment graph \( G = (V,E) \), whose vertices are locations,
and the dynamic interaction graph \( G^{A}(t) \), whose vertices are agents.

\subsection{Observations, policy and reward}

The multi-agent interaction is modelled as a partially observable stochastic
game
\begin{equation}
  \big( A,\ S,\ \{\mathcal{U}_i\}_{i},\ \{\mathcal{O}_i\}_{i},\ P,\ \{R_i\}_{i},\ \gamma \big),
\end{equation}
where \( S \) is the global state (agent positions, task set, storage
assignment) and \( P \) the transition kernel induced by the joint action and
the collision-resolution rule. The observation of agent \( a_i \) is
\begin{equation}
  o_i(t) = \Big(
    G_i(t),\ \eta_i,\
    \{ x_{i'}(t) : (a_i, a_{i'}) \in E^{A}(t) \},\
    \iota_i(t)
  \Big),
  \label{eq:observation}
\end{equation}
where \( x_{i'}(t) \) are features of neighbouring agents and \( \iota_i(t) \)
encodes task information (current pickup and delivery targets, and, for
section-based variants, the assigned zone).

Features are embedded by a graph neural network over \( G_i(t) \) and
\( G^{A}(t) \). With \( h_i^{(0)}(t) = x_i(t) \), layer \( \kappa \) computes
\begin{equation}
  h_i^{(\kappa+1)}(t)
  = \varphi\Big(
      W^{(\kappa)} h_i^{(\kappa)}(t)
      + \!\!\sum_{(a_i,a_{i'}) \in E^{A}(t)}\!\! W_{\mathrm{msg}}^{(\kappa)} h_{i'}^{(\kappa)}(t)
    \Big),
  \label{eq:gnn}
\end{equation}
with nonlinearity \( \varphi \). The action is drawn from
\begin{equation}
  u_i(t) \sim \pi_\theta\big( \cdot \mid h_i^{(L)}(t) \big),
  \qquad u_i(t) \in \mathcal{U}\big(\ell_i(t)\big),
\end{equation}
where inadmissible actions are masked out.

The per-agent reward decomposes into a movement term, a task-completion
term, a collision term and a congestion term,
\begin{equation}
  R_i(t)
  = w_{\mathrm{mov}} R_i^{\mathrm{mov}}(t)
  + w_{\mathrm{tsk}} R_i^{\mathrm{tsk}}(t)
  + w_{\mathrm{col}} R_i^{\mathrm{col}}(t)
  + w_{\mathrm{cng}} R_i^{\mathrm{cng}}(t),
  \label{eq:reward}
\end{equation}
where the congestion term penalises entering locally crowded regions of
\( G \) and rewards leaving them, measured by the local occupancy
\begin{equation}
  \delta_i(t)
  = \frac{\big| \{ a_{i'} \in A : d_G(\ell_i(t), \ell_{i'}(t)) \leq r_{\mathrm{cng}} \} \big|}
         {\big| \mathcal{N}_{r_{\mathrm{cng}}}(\ell_i(t)) \big|},
  \qquad
  R_i^{\mathrm{cng}}(t) = -\,\mathrm{sgn}\big( \delta_i(t) - \delta_i(t-1) \big).
  \label{eq:congestion-reward}
\end{equation}
Setting \( w_{\mathrm{cng}} = 0 \) yields the congestion-blind ablation.
The learning objective is
\begin{equation}
  \max_{\theta}\ J(\pi_\theta)
  = \mathbb{E}\Big[ \textstyle\sum_{t=0}^{T} \gamma^t R_i(t) \Big].
\end{equation}

\subsection{Coupled system and evaluation}

The coupled system alternates two timescales. On the fast timescale, at every
timestep, tasks arrive, the routing controller acts, and agents move subject
to \eqref{eq:vertex-collision} and \eqref{eq:edge-collision}. On the slow
timescale, every \( \Delta \) timesteps, the storage assignment is revised by
\eqref{eq:storage-update}, which changes the pickup vertices of future tasks
and therefore the distribution the routing controller faces.

Flow concentration is quantified on the edge-load distribution
\( \{ \mu(e) \}_{e \in E} \). Writing \( \hat{\mu}(e) = \mu(e) / \sum_{e'} \mu(e') \),
we report the normalised entropy
\begin{equation}
  H = -\frac{1}{\log |E|} \sum_{e \in E} \hat{\mu}(e) \log \hat{\mu}(e),
  \label{eq:entropy}
\end{equation}
so that low \( H \) indicates commons-like concentration of traffic on few
edges. The experimental design varies (i) coupling, \( \Delta \) finite versus
\( \Delta = \infty \); (ii) controller class, \( \Pi_{\mathrm{cen}} \),
\( \Pi_{\mathrm{sec}} \), \( \Pi_{\mathrm{dec}} \); (iii) congestion sensitivity,
\( \beta \) and \( w_{\mathrm{cng}} \) zero or nonzero; (iv) local communication,
\( r = 0 \) versus \( r > 0 \); and (v) load, the number of agents \( m \) and
the task frequency \( f \). Reported measures are
\( \bar{\zeta} \), \( \Lambda \), \( \bar{c} \), \( \sum_e \omega(e) \), \( H \),
runtime per timestep, and their degradation as \( m \) and \( f \) increase.

% ---------------------------------------------------------------------
\section{Background to be Covered}
% ---------------------------------------------------------------------

The following material must be understood and presented before the
contribution can be stated precisely.

\subsection{Warehouse operations and storage assignment}
\begin{itemize}
  \item Warehouse order-picking systems: picker-to-goods versus
        goods-to-picker paradigms; the pick rate (order lines per hour) as
        the dominant industrial key performance indicator; zoning strategies
        (no zoning, progressive zoning) and their effect on picker
        utilisation and travel.
  \item The Storage Location Assignment Problem (SLAP): objectives (travel
        reduction, space utilisation, picking efficiency), constraints
        (capacity, class-based storage, item affinity), and standard
        heuristics such as demand-frequency (ABC/turnover) placement and
        proximity-to-depot placement.
  \item The Dynamic Storage Location Assignment Problem (DSLAP): why
        placement must be revised as demand and inventory change, and how
        deep reinforcement learning has been applied to it.
  \item Integrated storage-location and order-picking planning: existing
        attempts to couple placement with picking policy, and the gap that
        remains regarding routing-level congestion feedback.
  \item AGV-assisted order picking: exact routing algorithms for
        single-block parallel-aisle warehouses, order batching and batch
        sequencing, and handover-zone architectures with heterogeneous
        (human plus robotic) workers.
\end{itemize}

\subsection{Multi-agent path finding}
\begin{itemize}
  \item The classical MAPF formulation on a graph \( G = (V,E) \):
        vertex and edge (swap) collision constraints, objective functions
        (makespan versus sum-of-costs / flowtime), and complexity results
        (NP-hardness of optimal solving; inapproximability of makespan below
        a constant factor).
  \item Optimal centralised solvers: Conflict-Based Search and its
        improvements (ICBS), Increasing Cost Tree Search, M*, Enhanced
        Partial Expansion A*, Independence Detection with Operator
        Decomposition, and reductions to SAT / ILP / ASP.
  \item Bounded-suboptimal and anytime solvers: prioritised and
        hierarchical cooperative A*, Push-and-Swap / Push-and-Rotate, large
        neighbourhood search (MAPF-LNS, MAPF-LNS2), and LaCAM for
        large-scale instances.
  \item MAPF benchmarks and standard maps, together with the definitions,
        variants and reporting conventions used in the field.
  \item Generalisations relevant to warehouses: kinematic constraints,
        delay probabilities, combined target assignment and path finding,
        and MAPF with deadlines.
  \item MAPF on non-lattice graphs: why grid-specific machinery
        (convolutional encoders, Euclidean direction vectors, coordinate
        inputs) does not transfer, and how distance-to-goal node potentials
        replace direction on arbitrary graphs.
\end{itemize}

\subsection{Lifelong and online extensions: MAPD}
\begin{itemize}
  \item The Multi-Agent Pickup and Delivery problem: online arrival of
        tasks, lifelong operation, tasks as pickup-delivery vertex pairs,
        free versus occupied agents, and service time as the evaluation
        measure.
  \item Well-formed instances: endpoints, task endpoints and non-task
        endpoints; the sufficient condition for solvability; why agents must
        rest only where they cannot block others.
  \item Decoupled MAPD algorithms: Token Passing (TP) and Token Passing with
        Task Swaps (TPTS), their completeness guarantees on well-formed
        instances, and the communication-versus-effectiveness trade-off
        against a centralised baseline.
  \item The empirical trade-off surface: effectiveness ordering (centralised,
        TPTS, TP) versus runtime ordering (TP, TPTS, centralised), and the
        scalability limits that motivate decentralised learned controllers.
  \item Lifelong MAPF in large-scale warehouses with bounded re-planning,
        and cooperative MAPD variants requiring explicit multi-agent
        cooperation beyond collision avoidance.
  \item The relationship between task assignment and path finding: why MAPD
        is usually split into a task-assignment subproblem and a path-finding
        subproblem, and what is lost by that split.
\end{itemize}

\subsection{Reinforcement learning and multi-agent reinforcement learning}
\begin{itemize}
  \item Markov decision processes, partially observable MDPs, and partially
        observable stochastic games; value functions, policy gradients and
        the advantage function.
  \item Core algorithms needed for implementation: deep Q-networks with
        experience replay and target networks; asynchronous advantage
        actor-critic (A3C); advantage actor-critic (A2C) and proximal policy
        optimisation (PPO).
  \item Multi-agent reinforcement learning: independent learning and its
        non-stationarity problem; centralised training with decentralised
        execution; parameter sharing among homogeneous agents; shared
        experience actor-critic.
  \item Practical mechanisms: invalid action masking and its bias; curriculum
        learning over environment size and density; learning from
        demonstrations to seed high-reward regions; reward shaping and its
        pitfalls.
  \item Evaluation methodology for deep RL: seed variance, stratified
        bootstrap confidence intervals, and reporting standards.
\end{itemize}

\subsection{Graph neural networks for multi-agent systems}
\begin{itemize}
  \item Graph convolutional networks and message passing: node and edge
        attributes, graph filters, permutation equivariance, and the
        depth-versus-receptive-field trade-off.
  \item GNNs for decentralised multi-robot path planning: agents as nodes,
        proximity as edges, dynamically computed adjacency from the local
        field of view, and learned local communication.
  \item Crowd-aware decentralised MAPF: density-based reward terms, local
        communication via GNNs, and their measured contribution in ablation.
  \item Prioritised and implicit-priority communication learning as
        alternatives to symmetric message passing.
  \item Local observation design on graphs: fixed-depth subgraph extraction,
        the quality-versus-cost trade-off in the depth parameter, and
        generalisation to arbitrary graph sizes.
\end{itemize}

\subsection{Hierarchical and multi-level control}
\begin{itemize}
  \item Feudal reinforcement learning and Feudal Multi-Agent Hierarchies:
        manager agents producing goals for worker agents.
  \item Macro-actions under partial observability, and action-space
        decomposition by spatial zones.
  \item Hierarchical MARL for warehouse order picking: a manager assigning
        warehouse zones to heterogeneous workers, workers selecting item
        locations within the assigned zone, and a low-level controller
        (typically A*) executing primitive navigation.
  \item Industrial heuristic baselines that hierarchical controllers must
        beat: Follow-Me, Pick-Don't-Move, and closest-task assignment.
  \item Bridging centralised and decentralised path finding: hybrid
        architectures and the coordination-versus-scalability trade-off.
\end{itemize}

\subsection{Self-organisation, congestion and commons dynamics}
\begin{itemize}
  \item Sociotechnical systems framing of automated warehouses: humans and
        robots as co-workers on shared infrastructure.
  \item Emergence and self-organisation in multi-agent systems: how locally
        beneficial decisions can reinforce one another into either efficient
        global organisation or pathological concentration.
  \item Tragedy-of-the-commons dynamics in shared infrastructure, and MARL
        studies of commons problems; individual versus congestion-sensitive
        reward structures.
  \item Congestion measurement on graphs: edge loads, queue lengths, waiting
        times, deadlock frequency, and concentration indices (entropy, Gini)
        over the edge-load distribution.
  \item Related routing formulations for context: vehicle routing problems
        and their relation to warehouse-internal transport.
\end{itemize}
